% !TEX root = ../main.tex

% DO NOT TOUCH! unless absolutely necessary (consider updating the template then)

\documentclass[11pt, oneside]{book} % TODO add 'final' option once document is complete
\usepackage[a4paper,
            width=150mm,
            top=35mm,
            bottom=35mm,
            bindingoffset=10mm,
            marginparwidth=30mm,
            headheight=5mm,
            headsep=10mm,
            footskip=15mm,
            hcentering]{geometry}
\usepackage[backend=biber, sorting=nyt, citestyle=numeric-comp, bibstyle=ieee, maxbibnames=999]{biblatex}
\usepackage[ngerman, american]{babel}
\usepackage{setspace}
\usepackage{fancyhdr}
\usepackage[pdfusetitle, hidelinks]{hyperref}
\usepackage{mathtools}
\usepackage[warnings-off={mathtools-colon,mathtools-overbracket}]{unicode-math}
\usepackage{amsthm}
\usepackage{upgreek}
\usepackage{nicefrac}
\usepackage{algorithm}
\usepackage[beginComment=, beginLComment=, endLComment=]{algpseudocodex}
\usepackage{listings}
\usepackage{mleftright}
\usepackage{graphicx}
\usepackage{csquotes}
\usepackage{parskip}
\usepackage{titlesec}
\usepackage{xcolor}
\usepackage{lipsum}
\usepackage{pgfplots}
\usepackage{tikzpagenodes}
\usepackage{bookmark}
\usepackage{xurl}
\usepackage{etoolbox}
\usepackage{comment}
\usepackage{booktabs}
\usepackage{tabularx}
\usepackage{emptypage}
\usepackage{letltxmacro}
\usepackage[all]{nowidow}
\usepackage[final]{microtype}
\usepackage[nottoc]{tocbibind}
\usepackage[shortcuts]{extdash}
\usepackage[detect-all]{siunitx}
\usepackage[skip=5pt]{subcaption}
\usepackage[obeyFinal]{todonotes}
\usepackage[shortlabels]{enumitem}
\usepackage[nodayofweek]{datetime}
\usepackage{totcount}
\usepackage[figure,table,lstlisting,algorithm]{totalcount}
\usepackage[capitalize, nameinlink, noabbrev]{cleveref}
\usepackage[format=hang, font={color=darkgray, singlespacing}, skip=5pt]{caption}
\usepackage[lang-warn=false]{datatool-base}
\usepackage[style=alttree, acronym, symbols, nomain, nonumberlist, shortcuts=ac]{glossaries-extra}

% custom macros
% !TEX root = ../main.tex

% math typesetting utilities
\newcommand*{\given}{\;\middle|\;}
\newcommand*{\hateq}{\mathrel{\widehat{=}}}
\newcommand*{\tuple}[1]{\left\langle #1 \right\rangle}
\newcommand*{\set}[1]{\left\{ #1 \right\}}
\newcommand*{\multiset}[1]{\left\{\!\left\{ #1 \right\}\!\right\}}
\newcommand*{\evalat}[2][]{\left. #2 \right\rvert_{#1}}
\newcommand*{\bmat}[1]{\begin{bmatrix} #1 \end{bmatrix}}
\newcommand*{\sbmat}[1]{\begin{bsmallmatrix} #1 \end{bsmallmatrix}}
\newcommand*{\dbmat}[1]{\everymath{\displaystyle} \begin{bmatrix} #1 \end{bmatrix}}
\newcommand*{\abs}[1]{\left\lvert #1 \right\rvert}
\newcommand*{\norm}[1]{\left\lVert #1 \right\rVert}
\newcommand*{\bvec}[1]{{\symbfit{#1}}}

% math operators
\DeclareMathOperator*{\argmax}{arg\,max}
\DeclareMathOperator*{\argmin}{arg\,min}
\DeclareMathOperator*{\Exp}{\symbb{E}}
\DeclareMathOperator*{\Prob}{\symbb{P}}
\newcommand*{\expectation}[2][]{\Exp_{#1} \left[ #2 \right]}
\newcommand*{\probability}[1]{\Prob \left[ #1 \right]}

% environments
\newtheorem{theorem}{Theorem}
\newtheorem{problem}{Problem} \crefname{problem}{Problem}{Problems}
\newtheorem{hypothesis}{Hypothesis} \crefname{hypothesis}{Hypothesis}{Hypotheses}


% graphics directory
\graphicspath{{graphics/}}

% bibtex references
\addbibresource{resources/references.bib}

% acronym and symbol counters
\newtotcounter{acronyms}
\glsdefpostlink{acronym}{\glsxtrifwasfirstuse{\stepcounter{acronyms}}{}}
\newtotcounter{symbols}
\def\oldglsxtrnewsymbol{} \let\oldglsxtrnewsymbol=\glsxtrnewsymbol
\def\glsxtrnewsymbol{\stepcounter{symbols}\oldglsxtrnewsymbol}

% acronyms and symbols glossary
\makenoidxglossaries
\glspatchtabularx
\setabbreviationstyle[acronym]{long-short}
\loadglsentries{resources/acronyms.tex}
\setglossarypreamble[acronym]{
  \glssetwidest{RWTH\;\;\:}
  \setlength{\parskip}{0pt}
}
\glsxtrprovidestoragekey{group}{}{\glsgroup}
\newcommand*{\newglsgroup}[2]{\glsxtrnewsymbol[description={}, group=#1]{group-#1}{\textbf{#2}}}
\loadglsentries{resources/symbols.tex}
\setglossarypreamble[symbols]{
  \glssetwidest{\ensuremath{\argmax_a f(a)}\;\;\:}
  \setlength{\parskip}{0pt}
  \let\glsnamefont\textnormal
}

% fonts
\setmainfont{STIXTwoText}[
  Extension=.ttf,
  Path=./fonts/,
  UprightFont={*-Regular},
  BoldFont={*-Bold},
  ItalicFont={*-Italic},
  BoldItalicFont={*-BoldItalic}
]
\setmathfont{STIXTwoMath}[
  Extension=.ttf,
  Path=./fonts/
]
\setsansfont{NotoSans}[
  Extension=.ttf,
  Path=./fonts/,
  UprightFont={*-Regular},
  BoldFont={*-Bold},
  ItalicFont={*-Italic},
  BoldItalicFont={*-BoldItalic},
  Scale=MatchLowercase
]
\setmonofont{NotoSansMono}[
  Extension=.ttf,
  Path=./fonts/,
  UprightFont={*-Regular},
  BoldFont={*-Bold},
  ItalicFont={*-Regular},
  BoldItalicFont={*-Bold},
  ItalicFeatures=FakeSlant,
  Scale=MatchLowercase
]

% line spacing
\onehalfspacing
\setlist{nosep}

% page layout
\pagestyle{fancy}
\fancyhf{}
\makeatletter
\if@twoside
  \fancyhead[LE]{\nouppercase{%
    \if@mainmatter \thechapter\enskip \fi
    \leftmark
  }}
  \fancyhead[RO]{\nouppercase{%
    \if@mainmatter \thesection\enskip \fi
    \rightmark
  }}
  \fancyfoot[LE,RO]{\thepage}
\else
  \fancyhead[R]{\nouppercase{%
    \if@mainmatter \thechapter\enskip \fi
    \leftmark
  }}
  \fancyfoot[R]{\thepage}
\fi
\fancypagestyle{plain}{
  \fancyhf{}
  \fancyfoot[R]{\thepage}
  \renewcommand{\headrulewidth}{0pt}
}
\fancypagestyle{title}{
  \fancyhf{}
  \fancyfoot[R]{\place, \@date}
  \renewcommand{\headrulewidth}{0pt}
}
\makeatother
\renewcommand{\chaptermark}[1]{\markboth{#1}{}}
\renewcommand{\sectionmark}[1]{\markright{#1}}

% chapter layout
\titlespacing*{\chapter}{0pt}{0pt}{10mm}
\titleformat{\chapter}[hang]{\huge\bf}{\thechapter}{1em}{\huge}

% custom abstract environment
\newenvironment{abstract}{
  \chapter*{Abstract}
}{
  \glsresetall
}

% float placement
\makeatletter
\setlength{\@fptop}{0pt}
\setlength{\@fpsep}{17pt}
\makeatother
\raggedbottom

% todonotes layout
\reversemarginpar
\setuptodonotes{
  size=\footnotesize,
  backgroundcolor=white,
  linecolor=gray,
  bordercolor=gray,
  tickmarkheight=5pt
}
\let\TODO\todo

% pgfplot version specification
\pgfplotsset{compat=1.18}

% math parentheses spacing adjustments
\mleftright

% inline math line break prevention
\binoppenalty=10000
\relpenalty=10000

% remap (, ), [, and ] to \left(, \right), \left[, and \right]
\begingroup
  \catcode`(\active \xdef({\left\string(}
  \catcode`)\active \xdef){\right\string)}
  \catcode`[\active \xdef[{\left\string[}
  \catcode`]\active \xdef]{\right\string]}
\endgroup
\mathcode`(="8000 \mathcode`)="8000
\mathcode`[="8000 \mathcode`]="8000

% algorithm settings
\algrenewcommand\alglinenumber[1]{\color{gray}\footnotesize #1\enspace}

% listings settings
\lstset{
  basicstyle={\ttfamily\singlespacing},
  frame=lines,
  numbers=left,
  numberstyle={\color{gray}\ttfamily\footnotesize},
  xleftmargin=7mm,
  framexleftmargin=7mm,
  belowcaptionskip=\dimexpr-0.8244023083\baselineskip+8.5pt\relax,
  showstringspaces=false,
  commentstyle={\color{gray}\itshape}
}
\renewcommand{\lstlistlistingname}{List of Listings}

% Listing style for PDDL
\lstdefinestyle{PDDL}{
    language=Lisp, % Use Lisp for PDDL-like syntax, often works well for PDDL
    basicstyle=\ttfamily\footnotesize, % Font style and size
    keywordstyle=\color{blue}\bfseries, % Blue and bold keywords
    commentstyle=\color{green!60!black}\itshape, % Green and italic comments
    stringstyle=\color{purple}, % Purple strings
    showstringspaces=false, % Don't show spaces in strings
    breaklines=true, % Allow line breaks within code
    breakatwhitespace=true, % Break lines at whitespace
    tabsize=2, % Tab size
    captionpos=b, % Caption below the listing
    escapeinside={\%*}{*)} % Allows LaTeX commands inside listings
}

% conditional lists
\newcommand{\listacronyms}{
  \ifnum\totvalue{acronyms}>0
    \printnoidxglossary[type=acronym, title={List of Acronyms}]
  \fi
}
\newcommand{\listsymbols}{
  \ifnum\totvalue{symbols}>0
    \printunsrtglossary[type=symbols, title={List of Symbols}]
  \fi
}
\newcommand{\listfigures}{
  \iftotalfigures
    \listoffigures
  \fi
}
\newcommand{\listtables}{
  \iftotaltables
    \listoftables
  \fi
}
\newcommand{\listalgorithms}{
  \iftotalalgorithms
    \listofalgorithms\addcontentsline{toc}{chapter}{List of Algorithms}
  \fi
}
\newcommand{\listlistings}{
  \iftotallstlistings
    \lstlistoflistings\addcontentsline{toc}{chapter}{List of Listings}
  \fi
}
\newcommand{\listreferences}{
  \emergencystretch=1em
  \printbibliography[title={List of References}, heading=bibintoc]
}
\newcommand{\listtodos}{
  \ifnum\value{@todonotes@numberoftodonotes}>0
    \listoftodos[List of TODOs]
  \fi
}

% SI units settings
\robustify\bfseries
\sisetup{list-final-separator=\text{, and }}

% citation settings
\newcommand{\creflastconjunction}{, and }
